\section{Implémentation}
\providecommand{\fichier}[1]{\texttt{\detokenize{#1}}}

\subsection{Outils}
Le projet fonctionne sous Ubuntu 26.04 (WSL2 sous Windows). FreeFem++ (version~4.15) assure le maillage, la
discrétisation en éléments finis P1 et la résolution ; Python~3.14, avec NumPy, Pandas et Matplotlib, lit les
résultats, calcule les indicateurs et trace les figures ; Streamlit fournit l'interface web ; Git conserve
l'historique des versions et \LaTeX{} sert à rédiger ce rapport.

\subsection{Architecture}
FreeFem++ calcule, Python orchestre et visualise. Les deux échangent des fichiers texte : FreeFem++ exporte,
pour chaque scénario, la température en chaque nœud (colonnes $x$, $y$, $T$) et la connectivité des
triangles, ce qui préserve le trou de l'air lors du tracé ; Python relit ces fichiers. Pour lancer un
scénario, Python écrit un petit fichier de paramètres que le script FreeFem++ lit, appelle FreeFem++ et
recueille un fichier de résumé contenant les indicateurs.

\begin{table}[h]
\centering\small
\begin{tabular}{p{5.6cm}p{9.4cm}}
\toprule
Fichier & Rôle \\
\midrule
\fichier{01_mesh.edp}, \fichier{02_stationary.edp} & maillage d'une plaque et solveur stationnaire (cas test~1) \\
\fichier{tests/convergence_P1.edp} & étude de convergence sur une solution fabriquée \\
\fichier{02b_wall_stationary.edp} & mur de trois couches, convection de Robin, régime stationnaire \\
\fichier{03_transient.edp} & même mur en régime transitoire (Euler implicite) \\
\fichier{07_wall_daily_cycle.edp} & mur soumis à une journée type, régime périodique \\
\fichier{05_room2d.edp}, \fichier{06_room2d_param.edp} & coupe de pièce en 2D, régime stationnaire (version paramétrée) \\
\fichier{08_room2d_cycle.edp} & pièce 2D, cycle jour/nuit, air intérieur libre, soleil et ventilation \\
\fichier{run_scenarios.py}, \fichier{run_cycles.py} & lancement automatique des scénarios et tableaux comparatifs \\
\fichier{plot_*.py}, \fichier{plots.py} & figures \\
\fichier{app.py}, \fichier{ui_confort.py} & interface Streamlit \\
\bottomrule
\end{tabular}
\caption{Scripts principaux du dépôt.}
\label{tab:scripts}
\end{table}

\subsection{Maillage}
Le mur est maillé par une grille régulière de $32\times20$ mailles (1~cm dans l'épaisseur), soit 693 sommets et
1\,280 triangles. La pièce est maillée par une grille régulière de mailles de 2~cm ($232\times166$) dont on
retire la zone d'air ; il reste 11\,250 sommets et 21\,024 triangles. Toutes les épaisseurs de couche étant
des multiples de la taille de maille, les interfaces entre matériaux coïncident avec des arêtes du maillage.
Les bords du trou de l'air portent l'étiquette de frontière 5 (surfaces intérieures).

\subsection{Matériaux et paramètres}
Les propriétés $\lambda$ et $\rho c$ sont des fonctions constantes par triangle (éléments P0), évaluées au
centre de chaque triangle. Des contrôles d'aires vérifient le découpage en zones. Les scénarios sont décrits
dans des fichiers de configuration (\fichier{scenarios.json}, \fichier{cycle_scenarios.json}), de sorte que
changer une conductivité, une absorption solaire ou un débit de ventilation ne demande aucune modification
du code de calcul.

\subsection{Résolution}
Le système linéaire est résolu par une méthode directe (UMFPACK). Dans les simulations de cycle jour/nuit,
la matrice est identique à chaque pas de temps : elle est factorisée une seule fois, et seul le second
membre est recalculé, ce qui rend les 25 jours simulés par scénario réalisables en quelques minutes pour
les cinq scénarios.
